\chapter*{Anexo C: Desarrollo del proyecto por fases técnicas}
\addcontentsline{toc}{chapter}{Anexo C: Desarrollo del proyecto por fases técnicas}
\label{cap:anexo_desarrollo}

\setcounter{section}{0}
\renewcommand{\thesection}{C.\arabic{section}}
\setcounter{table}{0}
\renewcommand{\thetable}{C.\arabic{table}}

Los capítulos principales de esta memoria describen la arquitectura de Open
CVN y sus resultados de forma independiente del proceso concreto que llevó a
construirla. Este anexo es una excepción deliberada a esa regla: su
propósito específico es documentar la trazabilidad del desarrollo, por lo
que identifica de forma explícita las unidades de trabajo internas (issues y
correcciones puntuales) que el repositorio utilizó para planificar y
registrar cada fase. El desarrollo se organizó en tres fases técnicas
sucesivas, cada una cerrada con su propia batería de pruebas antes de
iniciar la siguiente.

\section{Fase 1: generación estructural y normalización}
\label{sec:anexo_desarrollo_fase1}

Esta fase construyó la base de todo el trabajo posterior: la traducción
reproducible del paquete oficial CVN a artefactos Python y la normalización
de sus metadatos dispersos, descritas en las
Secciones~\ref{sec:pipeline_generacion_estructural}
y~\ref{sec:pipeline_normalizacion}. La Tabla~\ref{tab:anexo_desarrollo_fase1}
resume sus unidades de trabajo.

\begin{table}[htbp]
  \centering
  \footnotesize
  \renewcommand{\arraystretch}{1.45}
  \caption{Unidades de trabajo de la Fase 1, agrupadas bajo el epic de automatización de la traducción CVN XML/XSD.}
  \label{tab:anexo_desarrollo_fase1}
  \begin{tabular}{p{0.1\textwidth} p{0.8\textwidth}}
    \hline
    \textbf{Issue} & \textbf{Aportación} \\
    \hline
    \#11 & Infraestructura de proyecto para la generación de código \\
    \#12 & Bindings Pydantic estructurales a partir de los esquemas XSD de CVN \\
    \#13 & Normalización de los metadatos del manual de especificación y del modelo en árbol \\
    \#14 & Definición de las reglas de mapeo semántico y de la política de resolución de conflictos \\
    \#15 & Implementación del generador de modelos de dominio Pydantic \\
    \#16 & Pruebas automatizadas del pipeline de generación \\
    \#17 & Documentación y automatización del flujo de trabajo completo \\
    \#25 & Pipeline de integración continua para la validación de contribuciones \\
    \hline
  \end{tabular}
\end{table}

\section{Fase 2: modelo conceptual, formato Open CVN y parser}
\label{sec:anexo_desarrollo_fase2}

Esta fase construyó, sobre la base anterior, el modelo conceptual agnóstico,
el formato de intercambio Open CVN JSON y el parser y validador públicos,
descritos en las Secciones~\ref{sec:pipeline_modelos_dominio}
y~\ref{sec:pipeline_json_schema} y en el
Capítulo~\ref{cap:formato_herramienta}. La
Tabla~\ref{tab:anexo_desarrollo_fase2} resume sus unidades de trabajo.

\begin{table}[htbp]
  \centering
  \footnotesize
  \renewcommand{\arraystretch}{1.45}
  \caption{Unidades de trabajo de la Fase 2, agrupadas bajo el epic de esquema agnóstico, JSON Schema y parser.}
  \label{tab:anexo_desarrollo_fase2}
  \begin{tabular}{p{0.1\textwidth} p{0.8\textwidth}}
    \hline
    \textbf{Issue} & \textbf{Aportación} \\
    \hline
    \#42 & Investigación de opciones para derivar diagramas UML de los modelos Pydantic \\
    \#43 & Definición de la capa de extracción del modelo conceptual agnóstico \\
    \#44 & Generación de diagramas UML o asimilables a UML \\
    \#45 & Generación de JSON Schema a partir de los modelos de dominio \\
    \#46 & Definición del formato canónico Open CVN JSON \\
    \#47 & Definición del contrato unificado de parser y validador \\
    \#48 & Extracción de XML embebido desde CVN PDF \\
    \#49 & Validación de la importación de XML y JSON \\
    \#50 & Pruebas y documentación del flujo de trabajo del parser \\
    \hline
  \end{tabular}
\end{table}

\section{Fase 3: aplicación de gestión local}
\label{sec:anexo_desarrollo_fase3}

Esta fase construyó la herramienta local de línea de órdenes que demuestra
el uso completo de Open CVN, descrita en la
Sección~\ref{sec:formato_herramienta_local}, incluyendo los mecanismos de
importación semántica y asistida por modelos de lenguaje que se evalúan con
garantías parciales en el Capítulo~\ref{cap:evaluacion}. La
Tabla~\ref{tab:anexo_desarrollo_fase3} resume sus unidades de trabajo.

\begin{table}[htbp]
  \centering
  \footnotesize
  \renewcommand{\arraystretch}{1.45}
  \caption{Unidades de trabajo de la Fase 3, agrupadas bajo el epic de la aplicación de gestión de currículos.}
  \label{tab:anexo_desarrollo_fase3}
  \begin{tabular}{p{0.1\textwidth} p{0.8\textwidth}}
    \hline
    \textbf{Issue} & \textbf{Aportación} \\
    \hline
    \#61 & Definición del alcance del MVP de la aplicación y de la interfaz de línea de órdenes \\
    \#62 & Implementación del almacenamiento local con SQLite \\
    \#63 & Implementación de versiones maestra y derivadas del currículo \\
    \#64 & Implementación del flujo de importación y exportación de Open CVN JSON \\
    \#65 & Implementación de la edición y selección de currículo en el MVP \\
    \#66 & Implementación de la exportación a LaTeX desde Open CVN \\
    \#67 & Implementación de la generación de PDF y de la entrega de la vista previa \\
    \#68 & Pruebas y documentación del MVP de la aplicación \\
    \#69 & Spike de importación de PDF asistida por modelos de lenguaje \\
    \#70 & Importación semántica de CVN XML a Open CVN JSON \\
    \#71 & Endurecimiento de limitaciones y documentación \\
    \hline
  \end{tabular}
\end{table}

\section{Correcciones de alcance}
\label{sec:anexo_desarrollo_hotfixes}

Además de las tres fases planificadas, el desarrollo registró ocho
correcciones puntuales de alcance, aplicadas cuando la implementación reveló
que una fuente auxiliar del paquete oficial CVN (tablas de referencia,
catálogos, convenciones del manual funcional) no encajaba con la
planificación inicial. Estas correcciones no añaden funcionalidad nueva:
ajustan decisiones previas a la vista de evidencia concreta encontrada
durante la implementación, y se documentan aquí porque forman parte
legítima de la trazabilidad del proceso. La
Tabla~\ref{tab:anexo_desarrollo_hotfixes} las resume.

\begin{table}[htbp]
  \centering
  \footnotesize
  \renewcommand{\arraystretch}{1.45}
  \caption{Correcciones de alcance aplicadas durante el desarrollo.}
  \label{tab:anexo_desarrollo_hotfixes}
  \begin{tabular}{p{0.14\textwidth} p{0.76\textwidth}}
    \hline
    \textbf{Corrección} & \textbf{Motivo} \\
    \hline
    \#1 & Actualización de la convención de registro del runner de generación \\
    \#2 & Alineación del punto de entrada humano del proyecto con el protocolo de actualización de documentación \\
    \#3 & Ampliación de la documentación del paquete fuente oficial de CVN \\
    \#4 & Corrección del alcance estructural para artefactos auxiliares del paquete fuente \\
    \#5 & Capa de resolución de normalización para fuentes de referencia auxiliares \\
    \#6 & Realineación de la planificación para la integración semántica de catálogos auxiliares \\
    \#7 & Evaluación dinámica de la elegibilidad de enumeración en tablas de referencia \\
    \#8 & Trazabilidad de tipos envoltorio en el traspaso de datos normalizados \\
    \hline
  \end{tabular}
\end{table}

La concentración de correcciones en torno a las fuentes auxiliares del
paquete oficial CVN (tablas de referencia y catálogos) no es casual: es la
misma tensión entre heterogeneidad documental y necesidad de trazabilidad
que motiva la política semántica conservadora descrita en la
Sección~\ref{sec:pipeline_politica_semantica}, y confirma en la práctica que
esa cautela de diseño estaba justificada.
